\subsection{A competing dyadic fixed component in the actual tensor quotient}

\begin{proposition}[P-F64: the single-component tensor route fails]
\label{prop:tensorcompetingcomponent}
Use $n\ge4$, $M=2^{n-1}$, the actual triple-rectangular
Kronecker quotient $X_n$, and its chain curve $Y_n$ from
P-F60.  The fixed scheme $X_n^{\boldsymbol\mu_M}$ has
a nonzero point outside $Y_n$.  In particular $Y_n$ is
not its unique positive-dimensional irreducible component.
This refutes the additional single-component hypothesis
considered in P-F60, not the existence of a global Hilbert
pole.  Comparing contributions of different components
remains necessary.
\end{proposition}

\begin{proof}
We use a known support and balanced array, rather than
claiming a new construction of that array.  The support
and squared amplitudes below are those of
Franks--Reichenbach, \emph{Barriers for recent methods in
geodesic optimization}, equation~(2.5) and Lemma~2.5,
printed pp.~14--15 in arXiv:2102.06652v3:
\url{https://arxiv.org/pdf/2102.06652}.
The authors explicitly attribute the array to Kravtsov,
Theorem~1 with $k=0$, and note the earlier
Kravtsov--Lukshin Theorem~9.  We verify the particular
balance and stabilizer consequences needed here.

Put
\[
 a_i=2^{-(n-i+1)},\qquad c_i=1-a_i\quad(2\le i\le n),
 \qquad w_0=2^{-(n-1)},
                                                        \tag{P-F64a}
\]
and define
\[
 K_n=\sqrt{w_0}\,e_1\otimes e_1\otimes e_1
 +\sum_{i=2}^n\left[
 \sqrt{a_i}\,(e_i\otimes e_1\otimes e_i
                    +e_i\otimes e_i\otimes e_1)
 +\sqrt{c_i}\,e_{i-1}\otimes e_i\otimes e_i\right].
                                                        \tag{P-F64b}
\]
All its displayed coefficients are nonzero.
No two distinct support triples agree in two positions.
For a fixed first index $i$ with $2\le i<n$, the last
two indices are $(1,i),(i,1),(i+1,i+1)$; they differ
in both entries.  At first index one they are $(1,1)$
and $(2,2)$, and at first index $n$ they are $(1,n)$
and $(n,1)$, with the same property.
For second index one the remaining pair is $(i,i)$,
$1\le i\le n$; for second index $j\ge2$ the two pairs
are $(j,1),(j-1,j)$, again differing in both entries.
The third-index assertion follows by symmetry in the
last two factors.  This proves the support assertion
for every choice of two positions.

Consequently all three reduced density matrices are
diagonal.  Their diagonal values are one: for the
first factor these identities are
\[
 w_0+c_2=1,\qquad 2a_i+c_{i+1}=1\ (2\le i<n),
                      \qquad 2a_n=1,
                                                        \tag{P-F64c}
\]
and for the second and third factors they are
$a_i+c_i=1$ for $i\ge2$ and
\[
              w_0+\sum_{i=2}^n a_i=1.
                                                        \tag{P-F64d}
\]
The latter is the finite geometric sum
$2^{-(n-1)}+2^{-(n-1)}+\cdots+2^{-1}=1$.
Exactly the moment-map calculation and Kempf--Ness
criterion used in P-F60 now show that the
$SL_n^3$ orbit of $K_n$ is closed and nonzero.
Its image $q_n\in X_n$ is therefore nonzero, by
uniqueness of the closed orbit in the fiber of the
affine reductive quotient.

Let $\zeta$ be a primitive $M$th root and put
\[
 b_1=0,\qquad b_i=2^{n-i}\ (2\le i\le n),\qquad
 A=\operatorname{diag}(\zeta^{b_i}),\quad B=A^{-1}.
                                                        \tag{P-F64e}
\]
The element $(A,B,B)$ fixes $K_n$.  Its weight on
$(1,1,1)$ is zero, on $(i,1,i)$ or $(i,i,1)$ it is
$b_i-b_i=0$, and on $(i-1,i,i)$ it is
$b_{i-1}-2b_i$.  This is zero as an integer for
$i\ge3$ and is $-M$ for $i=2$.
Since $\sum_i b_i=M-1$,
\[
       \chi(A,B,B)=(\det A)^{-1}
                   =\zeta^{-(M-1)}=\zeta.
                                                        \tag{P-F64f}
\]
The grading transformation formula (P-F60l) therefore
implies $q_n\in X_n^{\boldsymbol\mu_M}$.
We do not need to determine the entire stabilizer of
$K_n$ for the argument that follows.

We claim that $K_n$ and the symmetric chain tensor
$T_n$ of P-F60 belong to different $GL_n^3$ orbits.
By the full stabilizer computation (P-F60c), the
three components of every stabilizer element of $T_n$
are scalar multiples of the same matrix $D_\eta$.
Conjugating separately in the three factors preserves
the property that the three components have spectra
which differ only by a nonzero scalar.  If $K_n$
were in that orbit, its stabilizer element $(A,B,B)$
would have this property.  Thus the eigenvalue sets
of $A$ and $A^{-1}$ would differ by a scalar.

Those exponent sets modulo $M$ are
\[
        P=\{0,1,2,4,\ldots,M/2\},\qquad -P.
                                                        \tag{P-F64g}
\]
If $\zeta^P=z\zeta^{-P}$, the presence of $1$ in
$\zeta^{-P}$ implies $z\in\zeta^P$, so $z=\zeta^k$
for some $k$ and $P=k-P$ modulo $M$.
Since $0,1\in P$, this forces both $k$ and $k-1$
to lie in $P$.  With representatives in $[0,M-1]$,
$k=0$ is impossible because $M-1\notin P$, and a
power of two $k\ge4$ has odd predecessor at least
three, also not in $P$.  The only remaining cases
are $k=1,2$.  The former fails because $1-2=M-1$
modulo $M$, and the latter because $2-4=M-2$ modulo
$M$.  Both residues exceed $M/2$ when $M\ge8$.
This proves that the two spectra are not scalar
translates and establishes the claimed inequivalence.

For clarity, the nonzero points of $Y_n$ are exactly
the grading orbit of $p_n$.  By P-F60 its coordinate
ring is $\mathbb C[u^s:s\in S_n]$, where
$S_n\subset M\mathbb Z_{\ge0}$ and $S_n/M$ is a
cofinite numerical semigroup.  Write $v=u^M$ and
$S'=S_n/M$.  For any positive $h\in S'$, localization
at $v^h$ is precisely $\mathbb C[v,v^{-1}]$:
for all sufficiently large $k$ both $kh+1$ and
$kh-1$ lie in $S'$, giving
$v^{\pm1}=v^{kh\pm1}(v^h)^{-k}$.
If one positive monomial is nonzero at a point, all
are nonzero, since $(v^a)^h=(v^h)^a$.
Thus these localizations cover exactly the complement
of the vertex.  Grading acts there as $v\mapsto z^M v$,
which is transitive on $\mathbb C^*$.

If $q_n$ belonged to $Y_n$, it would consequently be
a grading translate of $p_n$.  Realize that grading
translate by a product-linear transformation with
the required determinant character, as in P-F60.
Both its translate of $T_n$ and $K_n$ have closed
$SL_n^3$ orbits.  Equality of quotient images would
make these closed orbits equal, contradicting the
proved $GL_n^3$ inequivalence.  Hence $q_n\notin Y_n$.

Finally the grading stabilizer of any nonvertex point
of a positively graded affine cone is finite: some
positive-degree homogeneous function has nonzero
value there, which restricts the stabilizer to a
finite group of roots of unity.  The grading orbit
of $q_n$ therefore has dimension one and lies wholly
in $X_n^{\boldsymbol\mu_M}$, since grading commutes
with that finite subgroup.  Its irreducible closure
is not contained in $Y_n$ and is contained in a
positive-dimensional irreducible component of the
fixed scheme.  As $Y_n$ itself is a one-dimensional
component by P-F60, the fixed scheme has another
positive-dimensional component.  None of these
arguments computes the total Hilbert residue or
proves cancellation of that residue.
\end{proof}
